\chapter{Introducción}
\label{ch:introduccion}

\section{Motivación}
\label{sec:motivacion}

El cáncer es un conjunto de enfermedades de las que se tiene registro desde el antiguo Egipto y la Grecia clásica~\cite{acs2026historia}, pero que aun con los avances modernos continúa representando uno de los mayores desafíos sanitarios a nivel global, alcanzando más de 20 millones de nuevos diagnósticos anuales~\cite{sung2026global}. Dentro de este panorama, los cánceres del tracto gastrointestinal representan alrededor del 36\,\% de la mortalidad oncológica mundial~\cite{marin2022expression}. En este grupo se encuentra el cáncer colorrectal, posicionado como el tercero de mayor incidencia y la segunda causa de muerte por cáncer en el mundo~\cite{sung2026global}. 

Si bien hoy en día existen tratamientos quimioterapéuticos estándar basados en fármacos como el 5-fluorouracilo, el oxaliplatino y el irinotecán, cerca del 50\,\% de los pacientes afectados por esta enfermedad en estadios avanzados desarrolla resistencia y sufre recaídas~\cite{abdallah2016, marin2022expression, qi2020long}. Estos fracasos en el tratamiento se han asociado con el desarrollo de \gls{mdr} adquirida, un proceso en el que cobran especial relevancia proteínas transportadoras de la familia de \gls{abc}, las cuales expulsan los medicamentos fuera de la célula tumoral y reducen su concentración intracelular~\cite{abdallah2016, marin2022expression, qi2020long}.
 
La \gls{mrp1}, perteneciente a la familia de transportadores \gls{abc}, se ha asociado con la resistencia a distintos agentes quimioterapéuticos en cáncer colorrectal~\cite{abdallah2016, cao2017role, qi2020long}. En este contexto, la inhibición de su actividad constituye una estrategia de interés para reducir el eflujo de fármacos y, potencialmente, mitigar la resistencia frente a los tratamientos convencionales~\cite{cao2017role}.

El descubrimiento de fármacos mediante métodos tradicionales de laboratorio es un proceso lento y costoso, condicionado por el tamaño del espacio químico estimado en más de $10^{60}$ moléculas posibles~\cite{das2025generative, luukkonen2023artificial}. Además, el diseño de fármacos es un problema intrínsecamente multiobjetivo, ya que una molécula viable no solo debe ser capaz de interactuar con la proteína objetivo con alta afinidad, sino también cumplir con perfiles farmacológicos adecuados y una sintetizabilidad viable~\cite{luukkonen2023artificial, zhang2026molecular}.

Para abordar esta dificultad, el \gls{cadd} y los modelos generativos, como los \glspl{vae}, permiten comprimir la información molecular en un espacio latente continuo desde el cual generar nuevas estructuras~\cite{das2025generative}. En este enfoque un codificador aprende a mapear representaciones moleculares discretas a vectores continuos, mientras que un decodificador realiza el proceso inverso, reconstruyendo o generando nuevas moléculas a partir de puntos de dicho espacio latente~\cite{das2025generative}. No obstante, muestrear este espacio de manera aleatoria no garantiza la obtención de compuestos con actividad biológica real ni con propiedades farmacológicas adecuadas~\cite{das2025generative}, por lo que resulta necesario el diseño de métodos que permitan explorar el espacio latente de manera efectiva.

La exploración del espacio dirigida a moléculas con características de interés puede ser vista como un problema de optimización combinatoria. Entre los enfoques para abordarlo, las metaheurísticas han cobrado gran interés debido a que pueden encontrar soluciones destacadas en espacios de búsqueda de elevada complejidad y tamaño, sin incrementar la complejidad computacional considerablemente~\cite{talbi2009metaheuristics}. Específicamente, las metaheurísticas multiobjetivo permiten optimizar simultáneamente distintas metas que suelen entrar en conflicto entre sí, como lograr una alta afinidad con la proteína, asegurar que la molécula sea fácil de sintetizar y mantener un buen perfil farmacológico, ofreciendo un conjunto de soluciones compromiso~\cite{luukkonen2023artificial, talbi2009metaheuristics}.

Esta tesina propone un esquema computacional de dos etapas. En la primera etapa, se entrena un modelo generativo \gls{vae} con bases de datos moleculares para crear un espacio latente continuo de compuestos químicos válidos. En la segunda etapa, se integran algoritmos metaheurísticos multiobjetivo para guiar la búsqueda dentro del espacio latente, evaluando iterativamente la aptitud de las moléculas generadas según su afinidad de unión y sus propiedades farmacológicas deseadas.

Bajo esta premisa, la presente tesina de grado se centra en el desarrollo y evaluación de un enfoque computacional que combina la capacidad generativa de los \glspl{vae} con la optimización guiada de algoritmos metaheurísticos multiobjetivo, con el propósito de descubrir moléculas candidatas viables y eficaces para inhibir la proteína \gls{mrp1}.

\section{Objetivos}
\label{sec:objetivos}

El objetivo general de este trabajo final de grado consiste en desarrollar y evaluar un enfoque computacional híbrido que combine modelos generativos con algoritmos metaheurísticos multiobjetivo, con el fin de obtener moléculas capaces de inhibir la resistencia a quimioterápicos en cáncer colorrectal. A partir de este objetivo general, se plantean los siguientes objetivos específicos:

\begin{itemize}
    \item Diseñar, implementar y evaluar modelos metaheurísticos multiobjetivo para optimizar moléculas generadas por un modelo generativo.
    \item Generar moléculas que tengan propiedades estables y que sean eficaces para inhibir la resistencia a quimioterápicos en cáncer colorrectal.
\end{itemize}

\section{Organización del Documento}
\label{sec:organizacion}

Esta tesina de grado se organiza en capítulos donde se establecen los fundamentos teóricos, se detallan los aspectos metodológicos y experimentales, y se presentan los resultados y las conclusiones alcanzadas. El contenido se estructura de la siguiente manera:

\begin{itemize}
    \item \textbf{Capítulo 1 (Introducción):} Expone la motivación, el problema a abordar, los objetivos del trabajo y la estructura general de la tesina.
    \item \textbf{Capítulo 2 (Marco Teórico):} Brinda los fundamentos sobre \gls{cadd}, representación molecular, métricas de evaluación, modelos generativos (\gls{vae}), optimización multiobjetivo y la proteína \gls{mrp1}, junto con los trabajos relacionados en el área.
    \item \textbf{Capítulo 3 (Metodología):} Describe el enfoque propuesto, detallando el ajuste de hiperparámetros mediante fármacos genéricos, el diseño de inhibidores específicos contra \gls{mrp1} y la configuración experimental utilizada.
    \item \textbf{Capítulo 4 (Resultados y Discusión):} Presenta los resultados obtenidos, analizando la calidad de las moléculas generadas, la comparación con el estado del arte y el estudio de los mejores compuestos candidatos.
    \item \textbf{Capítulo 5 (Conclusiones y Trabajo Futuro):} Resume las principales conclusiones del trabajo, las dificultades encontradas y las líneas de trabajo futuro.
\end{itemize}
